\documentclass[conference]{IEEEtran}
\IEEEoverridecommandlockouts
\usepackage{cite}
\usepackage{amsmath,amssymb}
\usepackage{graphicx}
\usepackage{booktabs}
\usepackage{url}
\def\BibTeX{{\rm B\kern-.05em{\sc i\kern-.025em b}\kern-.08em
    T\kern-.1667em\lower.7ex\hbox{E}\kern-.125emX}}

\begin{document}

\title{Activation Placement and Memory Contention in\\Microcontroller-Class Neural Inference}

\author{\IEEEauthorblockN{Alp O. Bolukbasi}
\IEEEauthorblockA{\textit{Department of Informatics} \\
\textit{Karlsruhe Institute of Technology}\\
Karlsruhe, Germany}}

\maketitle

\begin{abstract}
Placing a neural network's activation buffer in a different on-chip SRAM
region changes inference time on an STM32U585 by three cycles out of
320\,898 when no other bus master is active, and by up to 31\,189 cycles
when a DMA engine is running in the same region. This paper reports both
numbers from a single firmware image, using a cycle-accurate telemetry
path whose own overhead is 118 cycles per record and whose same-region
control spread is zero. The uncontended effect is 9.3 parts per million
and is not a useful scheduling signal. Under contention the effect is
9.7 percent, and it is asymmetric: traffic in the largest SRAM region
costs arenas in the other two regions roughly 4.7 percent each without
sharing a region with either, while traffic in those two regions costs
each other nothing measurable. A sweep of the competing master's working
set shows an effect whose sign depends on which region that master
occupies. The results argue that activation placement is worth treating
as a schedulable resource on this class of device only when the memory
system is shared, and they identify the bus topology of the largest
region as the mechanism that a placement policy would have to model.
\end{abstract}

\begin{IEEEkeywords}
embedded machine learning, memory interference, scratchpad allocation,
real-time systems, microcontrollers
\end{IEEEkeywords}

\section{Introduction}

Neural network inference on microcontrollers is usually described as a
compute problem and optimised as one. The activation arena, the scratch
buffer a converted network uses to hold intermediate tensors, is treated
as a size constraint: it either fits in on-chip SRAM or the network does
not deploy. Where in SRAM it fits is left to the linker.

That is a reasonable default on a device with one bus master. It is less
obviously reasonable on a modern microcontroller, where DMA engines,
radio stacks and sensor peripherals move data through the same bus matrix
as the core. If the placement of the activation arena interacts with that
traffic, then placement is a scheduling decision rather than a linker
detail, and a runtime could make it.

This paper measures whether that interaction exists and how large it is
on one concrete platform. The contribution is not a placement policy. It
is a measurement, taken carefully enough to be trusted in either
direction, of an effect that is usually assumed rather than quantified.

Three results follow. Without a competing bus master the effect is 9.3
parts per million, which is resolvable by the instrument and useless to a
scheduler. With a DMA engine active the effect reaches 9.7 percent. The
effect is strongly asymmetric between regions, and the asymmetry does not
follow from region size or from whether the arena and the competing
master share a region.

\section{Background and Related Work}

Scratchpad allocation for embedded systems is a mature problem. Banakar
et al.\ established the case for software-managed on-chip memory over
caches in embedded designs \cite{banakar2002}, and both Avissar et al.\
\cite{avissar2002} and Steinke et al.\ \cite{steinke2002} formulated the
assignment of program objects to scratchpad as an optimisation over a
static cost model. Those cost models assume a single requester. The
question here is what happens to such a model when a second master is
active, which is the condition a sensor-carrying microcontroller is
normally in.

Memory interference between concurrent requesters is well studied on
multicore platforms with shared DRAM. MemGuard regulates per-core memory
bandwidth to restore temporal isolation \cite{yun2013}, and PALLOC
allocates pages so that cores occupy different DRAM banks, which is
placement used for exactly the purpose examined here \cite{yun2014}. Kim
et al.\ bound the resulting interference delay analytically
\cite{kim2014}, and the predictable execution model of Pellizzoni et
al.\ restructures tasks into phases that do not contend
\cite{pellizzoni2011}. These works target processors with caches, DRAM
controllers and multiple cores. A Cortex-M33 microcontroller with
uncached on-chip SRAM has none of those, so the mechanism has to be
re-established rather than assumed.

On the inference side, TensorFlow Lite Micro documents the arena as the
central memory abstraction for microcontroller deployment
\cite{david2021}, and MLPerf Tiny standardises how such deployments are
benchmarked \cite{banbury2021}. MCUNet and its successor attack the same
memory limit from the model side, reducing peak activation footprint so
that networks fit at all \cite{lin2020,lin2021}. All of them treat the
arena as a quantity. None of them treat its address as a variable.

The name of the system reported here follows the real-time literature's
use of laxity, the slack between a task's remaining execution and its
deadline \cite{liu1973}, since the question being asked is whether
placement buys slack that a scheduler could spend.

\section{System Design}

\subsection{Platform}

Measurements were taken on a B-U585I-IOT02A board carrying an STM32U585
with a Cortex-M33 core at 160\,MHz, four flash wait states and the
regulator in its highest voltage scaling mode. The instruction cache is
enabled and instruction fetches come from flash. The data cache on this
part covers only the external memory aperture from
\texttt{0x60000000} to \texttt{0x9FFFFFFF}, so internal SRAM is uncached
and every access from the core to the activation arena is a bus
transaction that nothing absorbs. That property is what makes the device
suitable for this measurement: an effect at the bus matrix is not hidden
behind a cache.

The device exposes three contiguous SRAM regions in the system address
space: 192\,KiB at \texttt{0x20000000}, 64\,KiB at \texttt{0x20030000}
and 512\,KiB at \texttt{0x20040000}. A fourth region in the low-power
domain was excluded, since it is not comparable in intended use.

Application code runs under ThreadX. Inference uses a network converted
by ST Edge AI Core 4.0.1 from a public human activity recognition model
trained on the WISDM accelerometer dataset \cite{kwapisz2010}. The
network takes a window of 24 samples across three axes and emits four
class scores. Its activation arena is 2\,944 bytes, which is small
relative to every region under test and is a stated limit of this study.

\subsection{Instrumentation}

Inference is timed with the Data Watchpoint and Trace cycle counter. The
counter is 32 bits wide and wraps every 26.8 seconds at this clock, so
wrap is detected and flagged rather than assumed absent. Each completed
inference emits a 32-byte record carrying a sequence number, a release
timestamp, the execution interval, and identifiers for the arena region
and the competing master. Records accumulate in a single-producer
single-consumer ring and leave the device in framed batches protected by
a CRC.

The instrument was measured before it was used. A null probe that
exercises the record path without running inference costs 14 cycles to
read the ring and 118 cycles to push a record, which is 0.037 percent of
one inference. A sequence gap in the received stream is the only evidence
of a dropped record, so the sequence number advances on drop rather than
on transmission.

\subsection{Placement Mechanism}

The arena is a statically allocated buffer in each region under test,
with identical alignment, selected by pointer at run time rather than by
a linker script or a compile-time switch. The converted network accepts
its activation buffer through a run-time call, which is what makes the
runtime rather than the linker the owner of placement. Model inputs and
outputs are allocated inside the activation buffer, so moving the arena
moves the network's tensors with it.

Selecting at run time is not a convenience. A relink of identical sources
moved the observed median by 85 cycles in earlier work on this system,
which is an order of magnitude larger than the uncontended effect being
measured. Any design that rebuilds between regions cannot resolve that
effect. Every figure in Section V therefore comes from one firmware
image, identified by the hash of its flashed binary rather than of its
ELF container, since a clean rebuild reorders debug sections while
leaving the image identical.

\subsection{Competing Master}

The competing master is a GPDMA1 channel performing memory-to-memory
transfers on a circular linked list, with source and destination both in
the region under test, at high channel priority and with source and
destination on different matrix ports. It takes no interrupt, and its
completion count is polled outside the measured interval.

The choice of a DMA engine rather than a second thread is the central
methodological decision of this work. There is one core and no
simultaneous multithreading, so two threads never execute at the same
instant. A thread that preempts the inference thread contributes its own
cycles to the measured interval, producing a curve that rises smoothly
with aggressor intensity and that is scheduler accounting rather than
memory interference. A DMA engine is the only independent master
available, and it is the only configuration in which the bus arbiter has
a decision to make.

\section{Methodology}

Each measurement cell is one arena region crossed with one competing
master configuration. Three arena regions, three aggressor regions and an
aggressor-off column give twelve cells. Two further cells place the arena
in the first region under two different labels, so that the spread
between them is a same-region control and an estimate of the noise floor
that does not depend on any assumption about the instrument.

Cells are visited in randomised order across repetitions, so that thermal
or clock drift does not accumulate against one region. The system is
warmed before recording and the null probe runs before each session.
Medians and 99th percentiles are reported; a mean alone is not, since the
distributions are narrow and multimodal and a mean conceals both
properties.

The reported capture contains 9\,062 records over 180 seconds, between
217 and 225 samples per cell, with no dropped records, no sequence gaps,
no CRC rejections and one detected counter wrap. Raw captures record the
build configuration, the clock and voltage settings and the flashed image
hash alongside the samples, and analysis refuses a capture missing any of
them.

\section{Results}

\subsection{Placement without Contention}

Table~\ref{tab:matrix}, first column, gives the uncontended medians. The
three regions span three cycles out of 320\,898, which is 9.3 parts per
million. The same-region control pair returned 320\,898 and 320\,901,
identical to the corresponding uncontended figures, so the spread of the
instrument on the median is zero cycles at this sample size.

The instrument therefore resolves the uncontended effect, and the effect
is too small to act on. Reporting it as a null result would understate
what the measurement shows, and reporting it as a finding would overstate
what it is worth.

\begin{table}[t]
\caption{Median inference time in cycles, by arena region and competing
master region. Parenthesised values are the change against the same
arena's uncontended median.}
\label{tab:matrix}
\centering
\begin{tabular}{lrrrr}
\toprule
Arena & None & DMA in R1 & DMA in R2 & DMA in R3 \\
\midrule
R1 (192\,KiB) & 320\,898 & 337\,877 & 320\,927 & 335\,843 \\
              &          & (+16\,979) & (+29) & (+14\,945) \\
R2 (64\,KiB)  & 320\,901 & 320\,960 & 337\,845 & 335\,847 \\
              &          & (+59) & (+16\,944) & (+14\,946) \\
R3 (512\,KiB) & 320\,899 & 320\,967 & 320\,937 & 352\,088 \\
              &          & (+68) & (+38) & (+31\,189) \\
\bottomrule
\end{tabular}
\end{table}

\subsection{Placement under Contention}

The remaining columns of Table~\ref{tab:matrix} change the picture by
four orders of magnitude. Sharing a region with the competing master
costs 16\,979 cycles for the first region and 16\,944 for the second,
which is 5.3 percent. For the third region it costs 31\,189 cycles, or
9.7 percent, roughly twice as much.

Three observations follow, and only the first was expected.

Sharing a region with an active master is expensive, and the cost is
consistent between the first two regions to within 35 cycles.

The first and second regions do not interfere with each other. A master
in one costs an arena in the other 29 and 59 cycles respectively, which
is 0.01 percent and close to the resolution floor. For the purpose of
this workload they behave as independent slaves.

A master in the third region costs arenas in the other two 14\,945 and
14\,946 cycles, about 4.7 percent each, without sharing a region with
either. The reverse does not hold: masters in the first and second
regions cost an arena in the third 68 and 38 cycles. The asymmetry is
therefore a property of the third region rather than of the pairings.

\subsection{Working Set of the Competing Master}

Table~\ref{tab:sweep} sweeps the competing master's buffer over four
sizes. The effect is between 423 and 568 cycles across that range, well
above the zero-cycle control spread, and its sign depends on which region
the master occupies. A master in the first or second region does less
damage as its buffer grows. A master in the third region does more. When
both the arena and the master are in the third region the response is
almost flat.

An aggregate over all cells returns a change of a few hundred cycles in
the falling direction, which is the residue of two opposite trends and
does not describe either. Reporting the sweep per cell is therefore not a
presentational choice.

\begin{table}[t]
\caption{Median inference time in cycles against the competing master's
working set size.}
\label{tab:sweep}
\centering
\begin{tabular}{llrrrr}
\toprule
Arena & Master & 1\,KiB & 4\,KiB & 8\,KiB & 16\,KiB \\
\midrule
R1 & R1 & 338\,414 & 337\,981 & 337\,898 & 337\,877 \\
R1 & R3 & 335\,420 & 335\,746 & 335\,816 & 335\,843 \\
R2 & R2 & 338\,391 & 337\,951 & 337\,869 & 337\,845 \\
R2 & R3 & 335\,421 & 335\,747 & 335\,819 & 335\,847 \\
R3 & R3 & 351\,997 & 352\,052 & 352\,066 & 352\,088 \\
\bottomrule
\end{tabular}
\end{table}

\subsection{A Real Workload}

The synthetic master saturates the path it occupies, which makes it a
ceiling rather than an operating point. To place a realistic figure on
the same axis, the device also runs a sensing pipeline: an inertial
sensor sampled at 26\,Hz feeding the network's input window, an
environmental sensor, and a microphone captured through the digital
filter peripheral at a measured 16\,230 samples per second against the
16\,233 its clock chain predicts. The audio path writes roughly 32\,KiB
per second, three orders of magnitude below what the synthetic master
moves, and its buffer lands in the third region, which the previous
subsection identifies as the least convenient one.

This pipeline is always on and cannot currently be disabled at run time,
so its cost is measured as a shift in the uncontended baseline rather
than isolated from the inertial and environmental paths. It is reported
as the cost of the sensing workload as a whole, and separating the
contributions requires a run-time switch that this implementation does
not provide.

\section{Discussion and Limitations}

The practical reading of Table~\ref{tab:matrix} is that a placement
policy on this device would have almost nothing to optimise on an idle
memory system and a great deal to optimise on a busy one. That is a
narrower claim than the premise this work started from, and it is the
claim the data supports.

The asymmetry of the third region is the most useful result and the least
explained one. Three candidate mechanisms are consistent with the
measurements. The region may hang off a matrix port that also carries
traffic to the other two, so that saturating it blocks them. Its accesses
may cost roughly twice as much individually, which would explain both the
cross-region penalty and the near doubling of the same-region figure.
Or its arbitration may be shared at a point upstream of the slaves. The
transfer counter recorded with each sample advances at an identical rate
in all three regions, so it does not discriminate between these, and
resolving the question requires the bus matrix and SRAM sections of the
device reference manual rather than a further experiment.

Several limits sit next to the results rather than after them. The arena
is 2\,944 bytes and cannot be varied, since it is a property of the
converted network rather than a parameter, so the study says nothing
about how the effect scales with arena size. Measurements were taken at
optimisation level zero, which lengthens the inference and therefore
dilutes any fixed-cost component of the effect. Each region contributes
one arena address, so region and address are confounded, although a
second address within one region reproduced that region's figure. The
sampling rate of the inertial sensor is 26\,Hz against the 20\,Hz of the
dataset the network was trained on, which affects classification quality
and not timing. Finally, the competing master occupied one DMA channel
for most of the study and a different one after an audio driver claimed
the first, and channel index is assumed not to affect arbitration at
equal priority.

The network stack deserves a separate caution. With the radio driver's
threads running, inference costs about four percent more, but those
threads run above the inference thread in priority and preempt it, so
that figure mixes preemption with interference. The distinction this work
makes between a thread and a DMA engine applies to its own measurements
and would have to be applied again to separate those two components.

\section{Conclusion and Future Work}

Activation placement on this microcontroller is worth 9.3 parts per
million when the memory system is idle and up to 9.7 percent when it is
not. The instrument that produced both figures resolves three cycles in
320\,898 with a zero-cycle same-region control, which is what makes the
first number a measurement rather than an absence.

Two directions follow. The first is explanatory: identifying the topology
that makes the largest SRAM region penalise the other two without
sharing with them would turn a measured asymmetry into a cost model. The
second is operational: a runtime that placed the arena according to which
regions are currently busy would need exactly that cost model, and the
per-region figures here are the data such a policy would learn from. The
measurement path already carries the identifiers a policy would need, so
the remaining work is the policy rather than the plumbing.

\begin{thebibliography}{00}

\bibitem{liu1973} C. L. Liu and J. W. Layland, ``Scheduling algorithms
for multiprogramming in a hard-real-time environment,'' \emph{Journal of
the ACM}, vol. 20, no. 1, pp. 46--61, 1973.

\bibitem{banakar2002} R. Banakar, S. Steinke, B.-S. Lee, M. Balakrishnan,
and P. Marwedel, ``Scratchpad memory: a design alternative for cache
on-chip memory in embedded systems,'' in \emph{Proc. CODES}, 2002, pp.
73--78.

\bibitem{avissar2002} O. Avissar, R. Barua, and D. Stewart, ``An optimal
memory allocation scheme for scratch-pad-based embedded systems,''
\emph{ACM Trans. Embedded Computing Systems}, vol. 1, no. 1, pp. 6--26,
2002.

\bibitem{steinke2002} S. Steinke, L. Wehmeyer, B.-S. Lee, and P.
Marwedel, ``Assigning program and data objects to scratchpad for energy
reduction,'' in \emph{Proc. DATE}, 2002, pp. 409--415.

\bibitem{yun2013} H. Yun, G. Yao, R. Pellizzoni, M. Caccamo, and L. Sha,
``MemGuard: memory bandwidth reservation system for efficient performance
isolation in multi-core platforms,'' in \emph{Proc. RTAS}, 2013, pp.
55--64.

\bibitem{yun2014} H. Yun, R. Mancuso, Z.-P. Wu, and R. Pellizzoni,
``PALLOC: DRAM bank-aware memory allocator for performance isolation on
multicore platforms,'' in \emph{Proc. RTAS}, 2014, pp. 155--166.

\bibitem{kim2014} H. Kim, D. de Niz, B. Andersson, M. Klein, O.
Mutlu, and R. Rajkumar, ``Bounding memory interference delay in
COTS-based multi-core systems,'' in \emph{Proc. RTAS}, 2014, pp.
145--154.

\bibitem{pellizzoni2011} R. Pellizzoni, E. Betti, S. Bak, G. Yao, J.
Criswell, M. Caccamo, and R. Kegley, ``A predictable execution model for
COTS-based embedded systems,'' in \emph{Proc. RTAS}, 2011, pp. 269--279.

\bibitem{david2021} R. David et al., ``TensorFlow Lite Micro: embedded
machine learning for TinyML systems,'' in \emph{Proc. MLSys}, 2021, pp.
800--811.

\bibitem{banbury2021} C. Banbury et al., ``MLPerf Tiny benchmark,'' in
\emph{Proc. NeurIPS Datasets and Benchmarks Track}, 2021.

\bibitem{lin2020} J. Lin, W.-M. Chen, Y. Lin, J. Cohn, C. Gan, and S.
Han, ``MCUNet: tiny deep learning on IoT devices,'' in \emph{Proc.
NeurIPS}, 2020, pp. 11711--11722.

\bibitem{lin2021} J. Lin, W.-M. Chen, H. Cai, C. Gan, and S. Han,
``MCUNetV2: memory-efficient patch-based inference for tiny deep
learning,'' in \emph{Proc. NeurIPS}, 2021, pp. 2346--2358.

\bibitem{kwapisz2010} J. R. Kwapisz, G. M. Weiss, and S. A. Moore,
``Activity recognition using cell phone accelerometers,'' \emph{ACM
SIGKDD Explorations Newsletter}, vol. 12, no. 2, pp. 74--82, 2010.

\end{thebibliography}

\end{document}
